%% kapa: exam=1 round=33 year=2022 period=2 form=A subject=appraisal_law subject_ko=감정평가관계법규 seq=29 number=29
%% source: Q-Net 감정평가사 기출문제 (cst00309) artlSeq=5209997 file=raw/1008/5209997/2022년도 제33회 감정평가사 1차시험 2교시(감정평가 관계법규, 회계학).pdf
%% conversion: kordoc 4.24.51 -> tools/gen1.py (auto); review=proofread-sonnet; flags=none
\begin{question}{29}
공간정보의 구축 및 관리 등에 관한 법령상 토지소유자가 지적소관청에 토지의 합병을 신청할 수 없는 경우를 모두 고른 것은?
\begin{framedtext}
ㄱ. 합병하려는 토지의 지목이 서로 다른 경우 \par
ㄴ. 합병하려는 토지의 소유자별 공유지분이 다른 경우 \par
ㄷ. 합병하려는 토지가 구획정리를 시행하고 있는 지역의 토지와 그 지역 밖의 토지인 경우 \par
\end{framedtext}
\begin{choices}
  \item ㄱ
  \item ㄷ
  \item ㄱ, ㄴ
  \item ㄴ, ㄷ
  \item ㄱ, ㄴ, ㄷ
\end{choices}
\end{question}
\answer{5}
